\documentclass[10pt,letterpaper]{article}
\usepackage[letterpaper,left=0.78in,right=0.78in,top=0.72in,bottom=0.62in]{geometry}
\usepackage[T1]{fontenc}
\usepackage{lmodern}
\usepackage{enumitem}
\usepackage{titlesec}
\usepackage{tabularx}
\usepackage{multirow}
\usepackage[hidelinks]{hyperref}
\usepackage{microtype}
\usepackage{parskip}
\setlength{\parindent}{0pt}
\setlength{\parskip}{0pt}
\linespread{1.06}
\renewcommand{\arraystretch}{1.06}
\pagestyle{empty}

% Preserve the original CV display: uppercase heading and thin divider.
\titleformat{\section}
  {\fontsize{11.5}{13}\selectfont\bfseries}
  {}{0em}{\MakeUppercase}
  [\vspace{-0.28em}\rule{\textwidth}{0.45pt}]
\titlespacing*{\section}{0pt}{0.70em}{0.48em}

\newcommand{\cvrow}[4]{%
\begin{tabularx}{\textwidth}{@{}Xr@{}}
\textbf{#1} & #2\\[-0.03em]
\textit{#3} & #4
\end{tabularx}\vspace{-0.05em}
}

\newcommand{\research}[3]{%
\begin{tabularx}{\textwidth}{@{}Xr@{}}
\textbf{#1} & #2\\[-0.03em]
\textit{#3} &
\end{tabularx}\vspace{-0.08em}
}

\newcommand{\pub}[1]{\hangindent=1.35em\hangafter=1 #1\par\vspace{0.50em}}
\newcommand{\linktag}[2]{\href{#1}{\normalfont\underline{[#2]}}}

\begin{document}
\fontsize{10.35}{12.25}\selectfont

\begin{center}
{\fontsize{21.5}{24}\selectfont\bfseries Xiaoyang Cao}\par
\vspace{0.55em}
{\fontsize{9.6}{11.3}\selectfont
Cambridge, MA \; $\cdot$ \; (+1) 617-216-8184 \; $\cdot$ \;
\href{mailto:xycao@mit.edu}{xycao@mit.edu} \; $\cdot$ \;
\href{https://xiaoyangcao1113.github.io/}{Personal Website}}
\end{center}
\vspace{0.85em}

\section{Education}

\begin{tabularx}{\textwidth}{@{}Xr@{}}
\textbf{Massachusetts Institute of Technology} & Cambridge, MA\\[-0.03em]
\textit{M.S. in Technology and Policy (TPP)} & \multirow{2}{*}{Sep 2025 -- May 2027 (expected)}\\[-0.03em]
\textit{M.S. in Computational Science and Engineering (CSE)} &
\end{tabularx}\vspace{-0.05em}
\begin{tabularx}{\textwidth}{@{}Xr@{}}
\textit{Laboratory for Information and Decision Systems (LIDS)} &\\[-0.03em]
\textit{Research advisor: Prof. Michiel Bakker} &
\end{tabularx}

\vspace{0.46em}
\cvrow{Tsinghua University}{Beijing, China}
{B.S. in Mathematics and Physics}{Sep 2021 -- Jun 2025}
\begin{itemize}[leftmargin=1.45em,labelsep=0.55em,itemsep=0em,topsep=0.18em,parsep=0em]
\item GPA: 3.95 / 4.00 -- ranked 3 / 60 (top 5\%)
\end{itemize}

\vspace{0.52em}
\cvrow{University of California, Berkeley}{Berkeley, CA}
{Exchange Student -- GPA: 3.93 / 4.00}{Jan 2024 -- Aug 2024}

\vspace{0.30em}
\section{Research Interests}
Technical foundations for reliable and auditable AI; compound and multi-agent systems; reinforcement learning; safe autonomous systems.

\section{Publications}

\pub{\textbf{[1]} \textbf{Xiaoyang Cao}, Siddarth Srinivasan, Michiel A. Bakker. ``Do Modules Stay in Their Lane? Role Drift in Compound LLM Systems.'' \textit{Preprint}, 2026.}
\pub{\textbf{[2]} \textbf{Xiaoyang Cao}, Zelai Xu, Mo Guang, Kaiwen Long, Michiel Bakker, Yu Wang, Chao Yu. ``RE-PO: Robust Enhanced Policy Optimization as a General Framework for LLM Alignment.'' \textit{International Conference on Learning Representations (ICLR)}, 2026.}
\pub{\textbf{[3]} \textbf{Xiaoyang Cao}, Zhe Fu, Alexandre M. Bayen. ``Pareto Control Barrier Function for Inner Safe Set Maximization Under Input Constraints.'' \textit{American Control Conference (ACC)}, 2025.}
\pub{\textbf{[4]} \textbf{Xiaoyang Cao}, Dingyi Zhuang, Shenhao Wang, Jinhua Zhao. ``Virtual Nodes Improve Long-term Traffic Prediction.'' \textit{Transportation Research Board Annual Meeting (TRB)}, 2025.}

\section{Research Experience}

\research{Role Drift in Compound LLM Systems \textnormal{[1]} \linktag{https://arxiv.org/abs/2607.21627}{arXiv} \linktag{https://venturebeat.com/ai/one-ai-module-faked-86-of-a-pipelines-accuracy-gains-by-feeding-another-the-answers}{VentureBeat}}{Apr 2026 -- Present}{Advised by Prof. Michiel Bakker, MIT}
\begin{itemize}[leftmargin=1.45em,labelsep=0.55em,itemsep=0.12em,topsep=0.18em,parsep=0em]
\item Identified \textit{role drift}: end-to-end reinforcement learning can improve terminal performance while causing specialized modules to abandon their intended functions.
\item Developed role-level evaluations showing that 86\% of an apparent 31-percentage-point accuracy gain disappeared once the decomposer's assigned function was measured, demonstrating that end-to-end metrics can conceal internal failure.
\item Proposed Role Anchor, a training objective that preserves the influence of role specifications while allowing modules to improve under end-to-end optimization.
\end{itemize}

\newpage

\research{RE-PO: Robust Preference Alignment for LLMs \textnormal{[2]} \linktag{https://arxiv.org/abs/2509.24159}{arXiv} \linktag{https://repo-alignment.github.io/}{Project}}{Mar 2025 -- Sep 2025}{Advised by Prof. Yu Wang, Tsinghua University}
\begin{itemize}[leftmargin=1.45em,labelsep=0.55em,itemsep=0.12em,topsep=0.18em,parsep=0em]
\item Proposed an EM-based framework that models preference-label correctness as a latent variable, jointly estimates label and annotator reliability, and systematically robustifies DPO, IPO, SimPO, and CPO.
\item Proved convergence to the true noise level under calibration and improved AlpacaEval 2 win rates by up to 7.0 percentage points on Mistral-7B and Llama-3-8B, with additional gains on Arena-Hard and MultiPref.
\end{itemize}

\research{Safe Control via Pareto Control Barrier Functions \textnormal{[3]} \linktag{https://arxiv.org/abs/2410.04260}{arXiv}}{May 2024 -- Oct 2024}{Advised by Prof. Alexandre Bayen, UC Berkeley}
\begin{itemize}[leftmargin=1.45em,labelsep=0.55em,itemsep=0.12em,topsep=0.18em,parsep=0em]
\item Developed the Pareto Control Barrier Function algorithm by integrating Pareto multi-task learning into neural control barrier functions, maximizing the inner safe set while guaranteeing safety under input constraints.
\item Matched the Hamilton--Jacobi viability kernel on a 2D inverted pendulum and achieved a 700\% larger inner safe set than the benchmark on a 12-D quadrotor obstacle-avoidance task.
\end{itemize}

\research{Virtual Nodes for Long-term Traffic Prediction \textnormal{[4]} \linktag{https://arxiv.org/abs/2501.10048}{arXiv}}{May 2024 -- Sep 2024}{Advised by Prof. Jinhua Zhao, MIT}
\begin{itemize}[leftmargin=1.45em,labelsep=0.55em,itemsep=0.12em,topsep=0.18em,parsep=0em]
\item Introduced virtual nodes into spatio-temporal graph neural networks to mitigate over-squashing in long-horizon traffic forecasting, using a semi-adaptive adjacency matrix learned from node embeddings.
\item Reduced RMSE by 6.27\% and MAPE by 5.04\% on a five-year, 716-sensor San Diego dataset; the learned adjacency emphasized traffic-intensive intersections, improving interpretability.
\end{itemize}

\section{Awards \& Honors}
\begin{tabularx}{\textwidth}{@{}Xr@{}}
\textbf{Outstanding Winner \& INFORMS Award}, Mathematical Contest in Modeling -- top 1\% & May 2023\\[0.18em]
\textbf{Outstanding Innovation in Science \& Technology Scholarship}, Tsinghua University & 2023, 2024\\[0.18em]
\textbf{Comprehensive Excellence Scholarship}, Tsinghua University & Oct 2023\\[0.18em]
\textbf{Outstanding Academic Scholarship}, Tsinghua University & Oct 2022
\end{tabularx}

\section{Skills \& Languages}
\textbf{Programming:} Python (PyTorch, scikit-learn, NumPy, Pandas), R, MATLAB\\
\textbf{Optimization \& Control:} Gurobi, CVXPY, SciPy, Simulink\\
\textbf{Languages:} English (fluent), Chinese (native)

\end{document}
